
\documentclass{article}
\usepackage{colortbl}
\usepackage{makecell}
\usepackage{multirow}
\usepackage{supertabular}

\begin{document}

\newcounter{utterance}

\centering \large Interaction Transcript for game `dond', experiment `semi\_en', episode 15 with Qwen3.5{-}4B.
\vspace{24pt}

{ \footnotesize  \setcounter{utterance}{1}
\setlength{\tabcolsep}{0pt}
\begin{supertabular}{c@{$\;$}|p{.15\linewidth}@{}p{.15\linewidth}p{.15\linewidth}p{.15\linewidth}p{.15\linewidth}p{.15\linewidth}}
    \# & $\;$A & \multicolumn{4}{c}{Game Master} & $\;\:$B\\
    \hline

    \theutterance \stepcounter{utterance}  
    & & \multicolumn{4}{p{0.6\linewidth}}{
        \cellcolor[rgb]{0.9,0.9,0.9}{
            \makecell[{{p{\linewidth}}}]{
                \texttt{\tiny{[P1$\langle$GM]}}
                \texttt{You are playing a negotiation game in which you have to agree on how to divide a set of items among you and another player.} \\
\\ 
\texttt{Rules:} \\
\texttt{(a) You and the other player are given a set of items. Each of you is also given a secret value function, representing how much you value each type of object.} \\
\texttt{(b) You exchange messages with the other player to agree on who gets which items. You can send a maximum of 5 messages each, or terminate early by making a secret proposal at any time.} \\
\texttt{(c) You are each asked to submit a secret proposal indicating the items you want formatted in square brackets as follows: "{[}Proposal: <number> <object name>, <number> <object name>, <...>{]}"} \\
\texttt{(d) If your proposals are complementary, i.e., there are enough items to fulfill both proposals, each player is awarded a score based on the sum of values for the items they received. Otherwise, both of you get zero points.} \\
\texttt{(e) Your goal is to maximize the score you receive.} \\
\\ 
\texttt{Let us start.} \\
\\ 
\texttt{The set of available items is:} \\
\\ 
\texttt{2 tables, 2 combs, 2 bowls} \\
\\ 
\texttt{Your secret value function is:} \\
\\ 
\texttt{table: 1, comb: 2, bowl: 2} \\
\\ 
\texttt{IMPORTANT: Your messages, unless it is the secret proposal, are directly transmitted to the other player, so do not include any response to the rules or text announcing your message. To make a secret proposal, use the indicated format. Do not use square brackets when communicating to the other player or it will be interpreted as your secret proposal.} \\
            }
        }
    }
    & \\ \\

    \theutterance \stepcounter{utterance}  
    & \multicolumn{4}{p{0.6\linewidth}}{
        \cellcolor[rgb]{0.9,0.9,0.9}{
            \makecell[{{p{\linewidth}}}]{
                \texttt{\tiny{[P1$\rangle$GM]}}
                \texttt{I will propose a division that maximizes my value while ensuring the other player also receives a positive value (to avoid the zero{-}point penalty).} \\
\\ 
\texttt{My valuation:} \\
\texttt{{-} Table: 1} \\
\texttt{{-} Comb: 2} \\
\texttt{{-} Bowl: 2} \\
\\ 
\texttt{Total items: 2 tables, 2 combs, 2 bowls.} \\
\texttt{Total value available: 2(1) + 2(2) + 2(2) = 8.} \\
\\ 
\texttt{If I take both combs (value 4) and one bowl (value 2), my total is 6.} \\
\texttt{The remaining items would be 1 table (value 1) and 1 bowl (value 2).} \\
\texttt{If the other player takes the remaining table and bowl, their value is 3.} \\
\texttt{Since 6 + 3 = 9? Wait, let's re{-}calculate total value.} \\
\texttt{Total items: 2 tables, 2 combs, 2 bowls.} \\
\texttt{My values: Table=1, Comb=2, Bowl=2.} \\
\texttt{Total sum of my values for all items = 2*1 + 2*2 + 2*2 = 2 + 4 + 4 = 10.} \\
\\ 
\texttt{Let's try to split the items to maximize my score while giving the other player enough to make the deal valid (complementary).} \\
\texttt{Option A: I take 2 combs and 1 bowl. Value = 2*2 + 2 = 6.} \\
\texttt{Remaining: 2 tables, 1 bowl.} \\
\texttt{If the other player takes 2 tables and 1 bowl, their value depends on their function (unknown). But if they value tables and bowls positively, they might get a good deal.} \\
\texttt{However, to be safe and ensure a win{-}win, I should offer the other player items I value less, or items they likely value highly. Since I don't know their values, I should aim for a split where I get the highest value items I can, and the remainder is substantial enough that they won't reject it (assuming they are rational and want to avoid 0).} \\
\\ 
\texttt{Actually, the rule says "If your proposals are complementary... Otherwise, both of you get zero points." This implies if we don't agree on a split that covers all items, we get 0. So I must propose a set of items that, combined with the other player's potential proposal, covers the whole set.} \\
\\ 
\texttt{Let's propose: 2 combs, 1 bowl.} \\
\texttt{Value to me: 2*2 + 2 = 6.} \\
\texttt{Remaining: 2 tables, 1 bowl.} \\
\texttt{If the other player takes the remaining 2 tables and 1 bowl, the total items are covered.} \\
\texttt{Is there a better split for me?} \\
\texttt{If I take 2 combs and 2 bowls: Value = 2*2 + 2*2 = 8} \\
            }
        }
    }
    & & \\ \\

    \theutterance \stepcounter{utterance}  
    & & & \multicolumn{2}{p{0.3\linewidth}}{
        \cellcolor[rgb]{0.9,0.9,0.9}{
            \makecell[{{p{\linewidth}}}]{
                \texttt{\tiny{[GM$|$GM]}}
                \texttt{continue} \\
            }
        }
    }
    & & \\ \\

    \theutterance \stepcounter{utterance}  
    & & & \multicolumn{2}{p{0.3\linewidth}}{
        \cellcolor[rgb]{0.9,0.9,0.9}{
            \makecell[{{p{\linewidth}}}]{
                \texttt{\tiny{[GM$|$GM]}}
                \texttt{I will propose a division that maximizes my value while ensuring the other player also receives a positive value (to avoid the zero{-}point penalty).} \\
\\ 
\texttt{My valuation:} \\
\texttt{{-} Table: 1} \\
\texttt{{-} Comb: 2} \\
\texttt{{-} Bowl: 2} \\
\\ 
\texttt{Total items: 2 tables, 2 combs, 2 bowls.} \\
\texttt{Total value available: 2(1) + 2(2) + 2(2) = 8.} \\
\\ 
\texttt{If I take both combs (value 4) and one bowl (value 2), my total is 6.} \\
\texttt{The remaining items would be 1 table (value 1) and 1 bowl (value 2).} \\
\texttt{If the other player takes the remaining table and bowl, their value is 3.} \\
\texttt{Since 6 + 3 = 9? Wait, let's re{-}calculate total value.} \\
\texttt{Total items: 2 tables, 2 combs, 2 bowls.} \\
\texttt{My values: Table=1, Comb=2, Bowl=2.} \\
\texttt{Total sum of my values for all items = 2*1 + 2*2 + 2*2 = 2 + 4 + 4 = 10.} \\
\\ 
\texttt{Let's try to split the items to maximize my score while giving the other player enough to make the deal valid (complementary).} \\
\texttt{Option A: I take 2 combs and 1 bowl. Value = 2*2 + 2 = 6.} \\
\texttt{Remaining: 2 tables, 1 bowl.} \\
\texttt{If the other player takes 2 tables and 1 bowl, their value depends on their function (unknown). But if they value tables and bowls positively, they might get a good deal.} \\
\texttt{However, to be safe and ensure a win{-}win, I should offer the other player items I value less, or items they likely value highly. Since I don't know their values, I should aim for a split where I get the highest value items I can, and the remainder is substantial enough that they won't reject it (assuming they are rational and want to avoid 0).} \\
\\ 
\texttt{Actually, the rule says "If your proposals are complementary... Otherwise, both of you get zero points." This implies if we don't agree on a split that covers all items, we get 0. So I must propose a set of items that, combined with the other player's potential proposal, covers the whole set.} \\
\\ 
\texttt{Let's propose: 2 combs, 1 bowl.} \\
\texttt{Value to me: 2*2 + 2 = 6.} \\
\texttt{Remaining: 2 tables, 1 bowl.} \\
\texttt{If the other player takes the remaining 2 tables and 1 bowl, the total items are covered.} \\
\texttt{Is there a better split for me?} \\
\texttt{If I take 2 combs and 2 bowls: Value = 2*2 + 2*2 = 8} \\
            }
        }
    }
    & & \\ \\

    \theutterance \stepcounter{utterance}  
    & & \multicolumn{4}{p{0.6\linewidth}}{
        \cellcolor[rgb]{0.9,0.9,0.9}{
            \makecell[{{p{\linewidth}}}]{
                \texttt{\tiny{[GM$\rangle$P2]}}
                \texttt{You are playing a negotiation game in which you have to agree on how to divide a set of items among you and another player.} \\
\\ 
\texttt{Rules:} \\
\texttt{(a) You and the other player are given a set of items. Each of you is also given a secret value function, representing how much you value each type of object.} \\
\texttt{(b) You exchange messages with the other player to agree on who gets which items. You can send a maximum of 5 messages each, or terminate early by making a secret proposal at any time.} \\
\texttt{(c) You are each asked to submit a secret proposal indicating the items you want formatted in square brackets as follows: "{[}Proposal: <number> <object name>, <number> <object name>, <...>{]}"} \\
\texttt{(d) If your proposals are complementary, i.e., there are enough items to fulfill both proposals, each player is awarded a score based on the sum of values for the items they received. Otherwise, both of you get zero points.} \\
\texttt{(e) Your goal is to maximize the score you receive.} \\
\\ 
\texttt{Let us start.} \\
\\ 
\texttt{The set of available items is:} \\
\\ 
\texttt{2 tables, 2 combs, 2 bowls} \\
\\ 
\texttt{Your secret value function is:} \\
\\ 
\texttt{table: 0, comb: 1, bowl: 4} \\
\\ 
\texttt{IMPORTANT: Your messages, unless it is the secret proposal, are directly transmitted to the other player, so do not include any response to the rules or text announcing your message. To make a secret proposal, use the indicated format. Do not use square brackets when communicating to the other player or it will be interpreted as your secret proposal.} \\
\\ 
\\ 
\texttt{I will propose a division that maximizes my value while ensuring the other player also receives a positive value (to avoid the zero{-}point penalty).} \\
\\ 
\texttt{My valuation:} \\
\texttt{{-} Table: 1} \\
\texttt{{-} Comb: 2} \\
\texttt{{-} Bowl: 2} \\
\\ 
\texttt{Total items: 2 tables, 2 combs, 2 bowls.} \\
\texttt{Total value available: 2(1) + 2(2) + 2(2) = 8.} \\
\\ 
\texttt{If I take both combs (value 4) and one bowl (value 2), my total is 6.} \\
\texttt{The remaining items would be 1 table (value 1) and 1 bowl (value 2).} \\
\texttt{If the other player takes the remaining table and bowl, their value is 3.} \\
\texttt{Since 6 + 3 = 9? Wait, let's re{-}calculate total value.} \\
\texttt{Total items: 2 tables, 2 combs, 2 bowls.} \\
\texttt{My values: Table=1, Comb=2, Bowl=2.} \\
\texttt{Total sum of my values for all items = 2*1 + 2*2 + 2*2 = 2 + 4 + 4 = 10.} \\
\\ 
\texttt{Let's try to split the items to maximize my score while giving the other player enough to make the deal valid (complementary).} \\
\texttt{Option A: I take 2 combs and 1 bowl. Value = 2*2 + 2 = 6.} \\
\texttt{Remaining: 2 tables, 1 bowl.} \\
\texttt{If the other player takes 2 tables and 1 bowl, their value depends on their function (unknown). But if they value tables and bowls positively, they might get a good deal.} \\
\texttt{However, to be safe and ensure a win{-}win, I should offer the other player items I value less, or items they likely value highly. Since I don't know their values, I should aim for a split where I get the highest value items I can, and the remainder is substantial enough that they won't reject it (assuming they are rational and want to avoid 0).} \\
\\ 
\texttt{Actually, the rule says "If your proposals are complementary... Otherwise, both of you get zero points." This implies if we don't agree on a split that covers all items, we get 0. So I must propose a set of items that, combined with the other player's potential proposal, covers the whole set.} \\
\\ 
\texttt{Let's propose: 2 combs, 1 bowl.} \\
\texttt{Value to me: 2*2 + 2 = 6.} \\
\texttt{Remaining: 2 tables, 1 bowl.} \\
\texttt{If the other player takes the remaining 2 tables and 1 bowl, the total items are covered.} \\
\texttt{Is there a better split for me?} \\
\texttt{If I take 2 combs and 2 bowls: Value = 2*2 + 2*2 = 8} \\
            }
        }
    }
    & \\ \\

    \theutterance \stepcounter{utterance}  
    & & & \multicolumn{4}{p{0.6\linewidth}}{
        \cellcolor[rgb]{0.9,0.9,0.9}{
            \makecell[{{p{\linewidth}}}]{
                \texttt{\tiny{[GM$\langle$P2]}}
                \texttt{{[}Proposal: 2 combs, 1 bowl{]}} \\
            }
        }
    }
     \\ \\

    \theutterance \stepcounter{utterance}  
    & & & \multicolumn{2}{p{0.3\linewidth}}{
        \cellcolor[rgb]{0.9,0.9,0.9}{
            \makecell[{{p{\linewidth}}}]{
                \texttt{\tiny{[GM$|$GM]}}
                \texttt{proposal} \\
            }
        }
    }
    & & \\ \\

    \theutterance \stepcounter{utterance}  
    & & & \multicolumn{2}{p{0.3\linewidth}}{
        \cellcolor[rgb]{0.9,0.9,0.9}{
            \makecell[{{p{\linewidth}}}]{
                \texttt{\tiny{[GM$|$GM]}}
                [0, 2, 1]
            }
        }
    }
    & & \\ \\

    \theutterance \stepcounter{utterance}  
    & & \multicolumn{4}{p{0.6\linewidth}}{
        \cellcolor[rgb]{0.9,0.9,0.9}{
            \makecell[{{p{\linewidth}}}]{
                \texttt{\tiny{[P1$\langle$GM]}}
                \texttt{The other player has submitted a secret proposal. It is now your turn to submit a proposal of your own.} \\
            }
        }
    }
    & \\ \\

    \theutterance \stepcounter{utterance}  
    & \multicolumn{4}{p{0.6\linewidth}}{
        \cellcolor[rgb]{0.9,0.9,0.9}{
            \makecell[{{p{\linewidth}}}]{
                \texttt{\tiny{[P1$\rangle$GM]}}
                \texttt{{[}Proposal: 2 combs, 1 bowl{]}} \\
            }
        }
    }
    & & \\ \\

    \theutterance \stepcounter{utterance}  
    & & & \multicolumn{2}{p{0.3\linewidth}}{
        \cellcolor[rgb]{0.9,0.9,0.9}{
            \makecell[{{p{\linewidth}}}]{
                \texttt{\tiny{[GM$|$GM]}}
                \texttt{proposal} \\
            }
        }
    }
    & & \\ \\

    \theutterance \stepcounter{utterance}  
    & & & \multicolumn{2}{p{0.3\linewidth}}{
        \cellcolor[rgb]{0.9,0.9,0.9}{
            \makecell[{{p{\linewidth}}}]{
                \texttt{\tiny{[GM$|$GM]}}
                [0, 2, 1]
            }
        }
    }
    & & \\ \\

    \theutterance \stepcounter{utterance}  
    & & & \multicolumn{2}{p{0.3\linewidth}}{
        \cellcolor[rgb]{0.9,0.9,0.9}{
            \makecell[{{p{\linewidth}}}]{
                \texttt{\tiny{[GM$|$GM]}}
                [[0, 2, 1], [0, 2, 1]]
            }
        }
    }
    & & \\ \\

\end{supertabular}
}

\end{document}
